\begin{bibelanalyse}{Josua}{3}{1-5}
    \analysistexte{
        \versmarker{1}{
            \hauptsatz{
                \konj{Da} \verbN{machte}\annot{Präteritum, Indikativ} sich \person{Josua} \konj{früh}\annot{Zeit} \verbN{auf},
            }

            \nebensatz[1]{
                \konj{und}\annot{Bindewort} sie \verbN{zogen}\annot{Präteritum, Indikativ} aus \ort{Sittim} 
                }

                \nebensatz[2]{
                    \konj{und}\annot{Bindewort} \verbN{kamen}\annot{Präteritum, Indikativ} an den \ort{Jordan}, 
                }

                \nebensatz[2]{
                    er \konj{und}\annot{Bindewort} alle \person{Kinder} \ort{Israels}; 
                }

                \nebensatz[2]{
                    \konj{und}\annot{Bindewort} sie \verbN{rasteten}\annot{Präteritum, Indikativ} \konj{dort}\annot{Ortsbezeichnung}, 
                }

                \nebensatz[2]{
                    \konj{ehe}\annot{Zeitwort} sie \verbN{hinüberzogen}\annot{Präteritum, Indikativ}.
                }
            }
        }
    
        

    
    \begin{hinweise}
        
        \item \vers{1}{Josua mach sich schon früh auf den Weg in Richtung Jordanfluss, nachdem die Kundschafter zurückkamen.}
        \item \vers{2}{Drei Tage lang lagerten sie an dem Fluss.}
        \item \vers{3-4}{Vorsteher gehen durch das Lager und informieren die Bevölkerung wie es weiter geht und wie sie sich verhalten sollen.}
        \item \vers{5}{Josua ermutigt das Volk.}
        
    \end{hinweise}

    \begin{beobachtungen}
        
        \item \vers{1}Da! Diese DA bezieht sich auf die vorgehenden Kapitel. Josua hat für die Vorbereitung Kundschafter in das Land geschickt um Jericho auszukundschaften. Als sie zurück kamen gaben sie Josua eine Positive Antwort. \bibleverse{Jos}(2:24). Die Kundschafter sind überzeugt, dass Gott ihnen das Land Übergeben wird. Josua zog mit dem Volk an den Jordan, der zu dieser Zeit über die  Ufer getreten ist.
        \item \vers{2-3}Drei Tage lagerten sie an dem Jordanfluss. Nach drei Tagen wurden die Meldung ausgegeben, das es los geht und wie sie sich verhalten sollen. Es gibt genau Anweisungen, zuerst die Bundeslade getragen von den Leviten (\bibleverse{VMos}(10:8)) . Das Volk wird nochmals gewarnt die Lade nicht zu berühren (\bibleverse{IVMos}(4:15)). Die Lade ist der Wegweiser. Alles Volk soll sich hinter der Lade befinden und sie als Wegweiser benutzen.
        \item \vers{4} Das Volk soll einen Abstand von der Lade halten, so dass nicht doch plötzlich diese überholt wird. In Vers \vers{3} steht das die Lade \betonung{des Herrn eures Gottes} gehört und so ist der Wegbreiter und Wegweiser der \herr{} selber. Wer den \herr N überholt kann so auf Abwegen geraten.
        \item \vers{5} Eine Ermutigung durch Josua. Sie gehen nicht einfach blinglings los, sonder wissen dass der Herr mit ihnen ist. Josua ermutigt darum sein Volk auf Gott zu vertrauen und sich darauf vorzubereiten. Er ermutigt sie indem er ihnen verspricht, dass Gott eingreifen und Wunder tun wird.
        
    \end{beobachtungen}
    \begin{aussageabsicht}      
        1
        & 
        Früh gehen alle Kinder Israels los 
        & 
        Sie gehen alle los um den Jordan zu überqueren um das verheissene Land einnehmne zu können.
        \\
        \hline
        2
        & 
        Drei Tage lagern sie am Jordan
        & 
        Vorbereitung des Volkes
        \\
        \hline
        3
        &
        Die Leviten sollen die BL tragen.
        &
        Schutz der Bevölkerung. Die Bundeslade darf nur von Leviten angefasst werden.
        \\
        \hline
        4
        &
        Abstand zwischen der Lade und dem Volk
        &
        Das Volk kennt den Weg nicht. Die Lade wird von Gott geführt. Nur er weiss den Weg.
        \\
        \hline
        5
        &
        Heiligt euch!
        &
        Gott wird Wunder tun, sie werden also Gott Hautnah erleben. Er wird mit ihnen im Lager sein.
        \\
        \hline        
    \end{aussageabsicht}
\end{bibelanalyse}

